\documentclass[10.5pt]{article}

% helpful packages
\usepackage{geometry}
\usepackage{indentfirst}
\usepackage{enumitem}
\usepackage{url}
\usepackage[urlcolor=blue, colorlinks=true, linkcolor=blue, citecolor=blue]{hyperref}
\usepackage{array}
\usepackage{bm}
\usepackage{float}
\usepackage{amsmath}
\usepackage[official]{eurosym}
\usepackage[T1]{fontenc}
\usepackage{textcomp}
\usepackage{longtable}
\usepackage{tabularx}  

% formatting tweaks
\pagestyle{empty}
\geometry{margin=0.8in}
\setlength\parindent{9pt}
\setlist{nosep}
\renewcommand\labelitemi{}
\setlength{\parindent}{0em}
\setlength{\parskip}{0em}

% new commands
\newcommand{\name}[1]{\begin{center}\section*{\LARGE #1}\end{center}}
\newcommand{\topinfo}[1]{\begin{center}\vspace{-0.2cm}#1\vspace{-0.2cm}\end{center}}
\newcommand{\resumesection}[1]{\vspace{0.05cm}\section*{\large #1}\vspace{-0.2cm}\hrule\vspace{0.25cm}}
\newcommand\tab[1][.37cm]{\hspace*{#1}}
\newcommand{\latex}{\LaTeX\xspace}

% Abstract indentation command
% Indents abstract text, renders in small font, with a subtle left rule
\newcommand{\paperabstract}[1]{%
  \vspace{.15cm}%
  \begin{list}{}{%
    \setlength{\leftmargin}{1.4em}%
    \setlength{\rightmargin}{0em}%
    \setlength{\topsep}{0pt}%
    \setlength{\partopsep}{0pt}%
    \setlength{\itemsep}{0pt}%
    \setlength{\parsep}{0pt}%
  }%
  \item[] {\small\textbf{\textit{Abstract:}} #1}%
  \end{list}%
  \vspace{.4cm}%
}

%\usepackage{palatino}
\usepackage{charter}
%\usepackage{libertine}

\begin{document}
%-------------------------------------------------------------------------------
% HEADER 
%-------------------------------------------------------------------------------

\begin{center}
{\LARGE \textbf{Amen Jalal}} \\[0.5em]
Website: \href{http://amenjalal.com}{amenjalal.com} \\
Email: u.jalal@lse.ac.uk \\
Address: 32 Lincoln's Inn Fields, \\
London WC2A 3PH, United Kingdom
\end{center}


\vspace{.2cm}
%-------------------------------------------------------------------------------
\resumesection{Employment}
%-------------------------------------------------------------------------------

\textbf{Harvard University}, Cambridge, US. 
\vspace{.1cm}
\begin{itemize}
	\item Assistant Professor, Harvard Kennedy School \hfill Starting August 2027
	\item Academy Scholar, Harvard Academy for International and Area Studies \hfill Starting August 2026
\end{itemize}



%-------------------------------------------------------------------------------
\resumesection{Education}
%-------------------------------------------------------------------------------
\textbf{The London School of Economics}, London, UK. 
\vspace{.1cm}
\hfill Expected July 2026
\begin{itemize}
	\item MRes/PhD in Economics
	\item Fields: Development; Labor; Environmental
  \item Advisors: Nava Ashraf, Oriana Bandiera, Gharad Bryan, Robin Burgess
\end{itemize}

\vspace{.3cm}
\textbf{Yale University}, New Haven, US. \hfill August 2013--May 2017 
\vspace{.1cm}
\begin{itemize}
	\item B.A. in Economics (with Distinction), \emph{cum laude}
	\item Ronald Meltzer Prize for an Outstanding Senior Essay in Economics 
\end{itemize}


\vspace{-.4cm}


%-------------------------------------------------------------------------------
\resumesection{Grants}
%-------------------------------------------------------------------------------
\textbf{Total funding as PI or co-PI: \$1,305,942} 
\setlength{\LTleft}{.67cm}
\begin{longtable}{@{}p{.9cm}p{2.9cm}p{12.2cm}@{}}
\hspace{-.7cm} \euro317,000      & G\textsuperscript{2}LM\textbar LIC (IZA) & Market Based Solutions to Gender Inequality in Pakistan's Labour Markets \\
\hspace{-.7cm} \pounds10,675 & IGC & Heat Insurance at Work \\
\hspace{-.7cm} \$21,540 & The Weiss Fund & Heat Insurance at Work \\
\hspace{-.7cm} \$264,987    & JPAL K-CAI & Heat Insurance at Work \\
\hspace{-.7cm} \pounds5,000 & LSE-HER & Rebuilding Lives, Not Just Homes: Addressing Trauma in Disaster Recovery \\
\hspace{-.7cm} \pounds190,000 & IGC & Immediate Relief, Delayed Recovery: Labor Market Impacts of Natural Disasters \\
\hspace{-.7cm} \pounds5,000   & LSE-STICERD & Immediate Relief, Delayed Recovery: Labor Market Impacts of Natural Disasters \\
\hspace{-.7cm} \pounds90,000  & CEPR-STEG & Immediate Relief, Delayed Recovery: Labor Market Impacts of Natural Disasters \\
\hspace{-.7cm} \$35,000       & Harvard University & Immediate Relief, Delayed Recovery: Labor Market Impacts of Natural Disasters \\
\hspace{-.7cm} \$160,586      & IFPRI & What Happens When Women Suddenly Stop Receiving Cash Transfers?  \\
\hspace{-.7cm} \$50,000       & IPA & What Happens When Women Suddenly Stop Receiving Cash Transfers?  \\
\hspace{-.7cm} \pounds19,500  & IGC & Equilibrium Effects of a Billion Trees: Evidence From Pakistan \\
\hspace{-.7cm} \pounds20,000  & LSE-RISF & The Illusion of Time: Gender Gaps in Job Search and Employment \\
\hspace{-.7cm} \pounds5,000   & LSE-STICERD & The Illusion of Time: Gender Gaps in Job Search and Employment \\
\end{longtable}
\vspace{-.4cm}

%-------------------------------------------------------------------------------
\resumesection{Presentations (* scheduled)}
%-------------------------------------------------------------------------------
\vspace*{-.4cm}
\setlength{\LTleft}{0pt}
\setlength{\LTright}{0pt}

\begin{longtable}{@{}p{2cm}@{\hspace{0em}}p{\dimexpr\linewidth-2cm-0.8em\relax}@{}}

\textbf{2026} &
University of British Columbia;
Oxford University; Paris School of Economics; Warwick University; Harvard Kennedy School;
Cornell University; Manchester University; LMU Munich; Harvard Women and Public Policy Program; Norwegian School of Economics; Uppsala University; University of Naples Federico II; Center for Global Development; World Bank SARGIL; NBER Summer Institute – Development*; IIES*; Stanford University*; Northwestern University*; University of San Diego*; Penn State University* \\
\noalign{\vskip 0.6em}

\textbf{2025} &
Northeastern Universities Development Consortium (NEUDC);
Stanford Institute for Theoretical Economics (SITE) – Gender;
IZA PhD Workshop in Labor and Behavioral Economics;
HEC Paris PhD Conference;
Warwick PhD Conference;
Discrimination and Diversity Workshop;
Royal Holloway University;
European Association of Young Economists (EAYE);
Finnish Centre of Excellence in Tax Systems Research;
University of Exeter;
Scotland and Northern England Conference in Applied Microeconomics;
The Canadian Labour Economics Forum Conference;
UC Berkeley DevPEC;
G\textsuperscript{2}LM\textbar LIC (IZA) Research Conference \\
\noalign{\vskip 0.6em}

\textbf{2024} &
LUMS University;
World Bank Inclusive Job Platforms Workshop \\
\noalign{\vskip 0.6em}

\textbf{2023} &
Paris Dauphine University; LUMS University;
LSE Environment Camp \\

\end{longtable}


\vspace{.3cm}


%-------------------------------------------------------------------------------
\resumesection{Working papers}
%-------------------------------------------------------------------------------
\textbf{{Screening Women Out? Pay Transparency in Job Postings}} (Job Market Paper) \vspace*{.2cm} \\
\emph{Selected for the ReStud and EALE 2026 job market tours. Outstanding Paper Award by the Discrimination and Diversity Workshop (2025) and European Economic Association.}

\paperabstract{Up to half of the gender pay gap stems from women's sorting into low-wage firms. Do women prefer amenities to wages, or face barriers to search? I tackle this question using data on 29 million job applications from Pakistan's largest job search platform, combined with firm and worker surveys and a field experiment mandating pay transparency. I document that large, high-paying firms are more likely to omit salaries in job ads and less likely to offer flexibility – an amenity women value slightly more than men. When pay is disclosed, men and women respond similarly to wages. But when undisclosed, behaviors diverge: men search randomly, while women sort negatively on pay. A theoretical framework shows that pay non-disclosure amplifies small gender differences in amenity preferences into large gender gaps in applications. To test whether transparency closes these gaps, I randomize mandatory versus optional pay disclosure in 20,088 jobs across 8,906 firms on the platform. The experiment leaves large-firm pay and amenities unchanged. Yet women's applications to these firms increase 95\%, and men's 59\%, reversing the gender gap in directed search.  This implies women do not knowingly ``buy'' flexibility with pay; rather, they turn to flexibility when its price is unknown. Meanwhile, large firms most exposed to mandated transparency become 30\% more likely to disclose pay post-experiment, suggesting they overestimated the costs of transparency.}


\textbf{{The Illusion of Time: Gender Gaps in Job Search and Employment}} \vspace*{.1cm} \\with Oriana Bandiera and Nina Roussille. \textbf{\textit{Revise \& Resubmit, Review of Economic Studies.}}

\paperabstract{In countries with low female employment, college-educated women often transition directly from education to homemaking. Does this reflect informed, forward-looking choices or unanticipated constraints? We study this question in Pakistan, where two-thirds of college-educated women remain out of the labor force. Tracking 2,400 college-graduating students, we document that men and women start their search with similar work aspirations, apply at similar rates, and receive comparable numbers of job offers. Yet a 27 pp employment gap emerges within six months post-graduation. This gap stems largely from timing: for women alone, there is a critical window, immediately post-graduation, during which job search is associated with much higher chances of employment. To test whether this relationship is causal, we randomize a modest incentive to apply early. By shifting search into the early window, the intervention raises women's employment by ~20\% but leaves men's employment unaffected, closing a third of the gender gap. Our evidence suggests that applying early enables women to start working before demands from the marriage market arise. Treatment effects are driven by women who underestimate how quickly these demands materialize, revealing an ``illusion of time.'' This illusion can be persistent since women in our sample recognize the barriers to employment faced by their female peers, but overestimate their own ability to overcome them.}


\textbf{{Can Market Competition Reduce Corruption?}} \vspace*{.1cm} \\with Muhammad Haseeb, Alex Quispe and Kate Vyborny

\paperabstract{Corruption remains a major obstacle to the delivery of public services in developing countries. We study whether competition between service delivery agents can mitigate corruption, leveraging exogenous changes to the market structure of agents responsible for delivering government cash transfers in Pakistan. A reform that increased the market power of these agents led to a 29.1 pp increase in the probability that a beneficiary had to pay an involuntary bribe to access the cash transfer. However, in areas with 1 standard deviation higher competition, this increase in bribe payments is almost completely eliminated. We rule out that mechanisms other than competition drive these results, such as strategic entry, changes in market access, and differences in monitoring efforts or cash recipient characteristics.}


\resumesection{Select works in progress}

\textbf{{Heat Insurance at Work}} \vspace*{.1cm} \\with Ashley Pople, Pol Simpson, Kartik Srivastava, Eddy Zou and Oriana Bandiera

\paperabstract{Heatwaves, intensified by climate change, hit the poorest the hardest. Many are exposed to dangerous temperatures through outdoor or factory-based work, and have limited access to adaptive resources. How can social protection evolve to address the growing losses caused by extreme heat? We conduct a randomized controlled trial among low-income workers in Delhi, India, comparing the effects of temperature-indexed heat alerts, anticipatory daily wage payments delivered before heat events, and same-day daily wage payments delivered during heatwaves. We measure impacts on labour supply, adaptive behaviours, health, and financial wellbeing. The design separately identifies the value of information, the value of cash, and the value of acting ahead of heatwaves.}


\textbf{{The Fundamental Need for Shelter: Evidence from the World’s Largest Home Reconstruction Program}}  \vspace*{.1cm} \\ with Canishk Naik and Pol Simpson

\paperabstract{Psychologists have long argued that human needs are met in order of priority: until basic needs such as shelter are secured, attention and effort remain focused on them, delaying higher-order pursuits. This idea has implications for a range of economic outcomes — shaping, for instance, how much background risk a household bears and, thus, how willing it is to make high-risk, high-reward investments in the future. Despite its influence, the idea that needs follow a hierarchy has rarely been tested using credible causal evidence on economic behaviour. We provide such a test using the reconstruction of flood-destroyed homes in Sindh, Pakistan. Exploiting staggered disbursements of home-reconstruction grants, we compare otherwise-similar households that received funding just before and just after a program expansion. Combining administrative records with a survey of 6,000 households, we estimate the effects of home security on risk-taking, migration, health, schooling, and productive investment.}


\textbf{{Direct Damage, Indirect Costs: The Incidence of Natural Disasters}} \vspace*{.1cm} \\with Pol Simpson 

\paperabstract{Disaster-prone communities have long developed strategies to weather seasonal shocks, including migration, borrowing, and insurance. But these mechanisms rely on the shocks being localised and predictable. We examine what happens when they are not, using panel data on 5,100 households affected by unprecedented and widespread floods in Pakistan, and a recentered instrumental variables design that distinguishes the effects of local from regional flooding. Households experiencing greater local flooding suffer more physical damage and displacement, yet appear to cope: they receive more aid, travel farther for work, and sustain consumption across both years. Conditional on own flooding, however, greater regional exposure tells a different story. These households receive less aid—consistent with zero-sum allocation of limited relief—face sharp declines in agricultural employment, and are less able to commute to less-affected areas. In the first year, they cope by selling assets. By the second year, these buffers appear exhausted: households draw down savings, take on new debt, skip meals, and report worsening physical health. While disaster response policy tends to focus on emergency relief or long-term recovery, our findings reveal substantial medium-term costs—driven not by direct damage but by indirect and dynamic region-wide effects.}


\resumesection{Publications}

\textbf{{What Happens When Women Suddenly Stop Receiving Cash Transfers?}} \vspace*{.1cm} \\
with Nasir Iqbal, Mahreen Mahmud, and Kate Vyborny. \textbf{\textit{Journal of Public Economics}}.

\paperabstract{Female-targeted cash transfers are widely used as a policy tool to enhance women's empowerment. However, little is known about what happens when payments stop – whether due to budget cuts, program changes, or recipient graduation. We study how women are affected by program exit, and whether they experience backlash when they stop bringing cash home. Using a regression discontinuity-in-differences design around a revised eligibility threshold, we follow a panel of 2,333 women exiting Pakistan's largest cash transfer program. One year after exit, drawing on a comprehensive battery of empowerment, intimate partner violence, and well-being measures, we find no evidence of backlash. These results suggest that the gains experienced by women during the program are not undone by adverse reactions upon exit.}

%-------------------------------------------------------------------------------
\resumesection{Policy Writing}
%-------------------------------------------------------------------------------

\href{https://documents1.worldbank.org/curated/en/099155004142238180/pdf/P1763410d1e1af00108e170e5754d04fed9.pdf}{Using Biometrics to Deliver Cash Payments to Women.}  World Bank 2022.

\href{https://documents1.worldbank.org/curated/en/099406503032591629/pdf/IDU157126b0b1079f1483d18ce317cfccf04f886.pdf}{Collecting Accurate Data on Intimate Partner Violence.}  World Bank 2025.

\href{https://www.theigc.org/blogs/billion-tree-tsunami-afforestation-pakistan}{Equilibrium Effects of a Billion Trees on Ecosystems and Livelihoods.} IGC 2025.

\href{https://steg.cepr.org/sites/default/files/2025-04/PB024%20Simpson%20%28LRG%201646%29.pdf}{How the Rural Poor Cope with a Climate Catastrophe: Evidence from Pakistan's 2022 Floods.} STEG 2025.

%-------------------------------------------------------------------------------
\resumesection{Professional Service}
%-------------------------------------------------------------------------------
\textbf{Refereeing:} American Economic Review; Journal of Economic Behavior \& Organization; Economic Development and Cultural Change; Economica

\textbf{Organizing Committee:} LSE Environment Week (2023-24); LSE Environment Camp (2023-26)

\textbf{Paper selection:} European Association of Young Economists (2025); Pathways to Development (2023)

\textbf{Mentoring:} LSE Application Mentoring Programme (2022, 2023)

%-------------------------------------------------------------------------------
\resumesection{Teaching}
%-------------------------------------------------------------------------------
\textbf{The London School of Economics}, London UK \hfill September 2021-2025\vspace*{.2cm} \\
Graduate Teaching Assistant: \vspace*{.1cm}\\
\begin{tabular}{p{0.23\linewidth} p{0.23\linewidth} p{0.6\linewidth}}
Microeconomics I & Department of Economics & Bachelors \\
Macroeconomics I & Department of Economics & Bachelors \\
Microeconomics   & School of Public Policy & Masters of Public Administration \\
Empirical Methods & School of Public Policy & Exec. Masters of Public Administration / Public Policy \\
Introduction to Statistics & School of Public Policy & Exec. Masters of Public Administration / Public Policy \\
\end{tabular}


\end{document}
